\documentclass[11pt]{article}
\usepackage[margin=1in]{geometry}
\usepackage[T1]{fontenc}
\usepackage{graphicx}
\usepackage{hyperref}
\usepackage{enumitem}
\setlength{\parskip}{0.5em}
\setlength{\parindent}{0pt}

\title{diffusion models for video generation}
\author{arXiv GraphRAG --- phase-6 deep research}
\date{\today}

\begin{document}
\maketitle

\begin{abstract}
A highly professional review that captures the major architectural shifts but could be strengthened by more explicit coverage of the Diffusion Transformer (DiT) paradigm and the emerging 'world model' perspective.
\end{abstract}

\section{From Pixels to Motion: A Comprehensive Review of Diffusion Models for Video Generation}

The transition from static image synthesis to the generation of coherent, dynamic video sequences represents one of the most significant challenges in modern artificial intelligence. While generative modeling has achieved remarkable success in the spatial domain, video generation introduces the complex dimension of time, requiring models to not only synthesize realistic textures but also maintain logical motion and structural integrity across sequential frames. The ability to generate high-fidelity video from minimal input is critical for industries ranging from cinematic entertainment to autonomous vehicle simulation, making the development of robust diffusion-based frameworks a central pursuit in the field.

This review traces the technological evolution of these systems, beginning with the Foundations of Generative Modeling and Diffusion to establish the mathematical principles of noise-based synthesis. We then examine the transition into the temporal domain through Early Video Generation and Latent Diffusion Models, which laid the groundwork for moving beyond single-frame generation. As the field matured, research shifted toward user-centric capabilities, specifically through Advancements in Text-to-Video and Motion Control, allowing for more precise semantic and physical guidance over the generated content.

To address the inherent instabilities of sequence generation, this review further investigates Temporal Consistency and Video Editing Techniques, which focus on mitigating flickering and enabling fine-grained manipulation of existing footage. We also explore the computational evolution of the field, looking at Efficient Architectures and Long-Term Generation to understand how models are being scaled for extended durations and practical deployment. Finally, the discussion culminates in Current Frontiers: Survey and Advanced Representations, providing a synthesis of the most recent breakthroughs and the emerging paradigms that define the current state of the art.

\subsection{Foundations of Generative Modeling and Diffusion}

The evolution of generative modeling has transitioned from variational frameworks and adversarial games to the current dominance of diffusion-based processes. Early foundational work established stochastic variational inference through a reparameterization of the variational lower bound to enable optimization via stochastic gradient methods \texttt{arXiv:1312.6114}. This was followed by the introduction of Generative Adversarial Networks (GANs), which shifted the paradigm toward a minimax two-player game between a generator and a discriminator \texttt{arXiv:1406.2661}. While GANs offered a novel adversarial framework, subsequent research sought to improve their stability, such as the Two Time-Scale Update Rule (TTUR) which utilizes individual learning rates to ensure convergence to a local Nash equilibrium \texttt{arXiv:1706.08500}. Further refinements in adversarial modeling included style-based architectures designed to achieve unsupervised separation of high-level attributes and scale-specific synthesis control \texttt{arXiv:1812.04948}.

A distinct lineage of generative modeling emerged from non-equilibrium thermodynamics, where an iterative forward diffusion process destroys data structure, which is then restored via a learned reverse process \texttt{arXiv:1503.03585}. This concept was further developed through score matching, which estimates gradients of perturbed data distributions across multiple noise levels using annealed Langevin dynamics \texttt{arXiv:1907.05600}. These disparate paths—variational bounds, adversarial games, and score-based gradients—converged into the modern diffusion paradigm. Denoising Diffusion Probabilistic Models (DDPMs) unified these ideas by training on a weighted variational bound connected to denoising score matching, achieving high-quality image synthesis with an Inception score of 9.46 on CIFAR-10 \texttt{arXiv:2006.11239}.

The efficiency and theoretical unity of diffusion models have been central themes in recent advancements. To address the slow sampling speeds inherent in Markovian diffusion, Denoising Diffusion Implicit Models (DDIMs) introduced non-Markovian processes that allow for 10x to 50x faster sampling without altering the training objective \texttt{arXiv:2010.02502}. Theoretical cohesion was further achieved through a stochastic differential equation (SDE) framework, which unifies score-based modeling and diffusion probabilistic modeling into a single predictor-corrector or neural ODE framework \texttt{arXiv:2011.13456}. This unification has allowed diffusion models to eventually surpass GANs in image synthesis quality; for instance, by employing classifier guidance to trade off diversity for fidelity, diffusion models achieved competitive FID scores on ImageNet (e.g., 3.85 on ImageNet 512x512 with upsampling and guidance) \texttt{arXiv:2105.05233}. These foundational architectures often leverage established components like the U-Net for precise localization through symmetric expanding paths \texttt{arXiv:1505.04597}, or the Transformer's attention mechanisms to dispense with recurrence and convolutions \texttt{arXiv:1706.03762}.

\subsection{Early Video Generation and Latent Diffusion Models}

The evolution of video generation has progressed from early adversarial frameworks focused on scene dynamics to sophisticated latent diffusion models capable of high-fidelity, text-conditioned synthesis. Initial methodologies primarily utilized Generative Adversarial Networks (GANs) to tackle the complexities of temporal modeling. For instance, \texttt{arXiv:1609.02612} introduced a spatio-temporal convolutional architecture designed to untangle foreground from background to learn scene dynamics. This approach was extended by \texttt{arXiv:1611.06624}, which proposed Temporal Generative Adversarial Nets (TGAN) using a dual-generator architecture and singular value clipping to stabilize training. Further refinement in decomposition was achieved by MoCoGAN (\texttt{arXiv:1707.04993}), which utilized an unsupervised adversarial learning scheme to separate motion and content via distinct latent vectors. While these early GAN-based models established foundational techniques for semantic representation, they faced challenges in scaling. This was addressed by the Dual Video Discriminator GAN (DVD-GAN) in \texttt{arXiv:1907.06571}, which employed a computationally efficient discriminator decomposition to enable the synthesis of longer, higher-resolution videos.

As the field shifted toward diffusion-based paradigms, the focus moved from purely adversarial training to modeling data distributions through denoising. The transition was catalyzed by advancements in image synthesis, such as the autoregressive transformer approach in \texttt{arXiv:2102.12092} and the comparative study of guidance techniques in GLIDE (\texttt{arXiv:2112.10741}). A pivotal breakthrough occurred with the introduction of Latent Diffusion Models (LDMs) in \texttt{arXiv:2112.10752}, which significantly reduced computational complexity by operating within the latent space of pretrained autoencoders. This efficiency allowed for general conditioning via cross-attention layers, a concept further enhanced by \texttt{arXiv:2205.11487}, which integrated large pretrained text-only transformers (e.g., T5) to improve image-text alignment and photorealism.

These image-centric advancements provided the architectural blueprint for video-specific diffusion. \texttt{arXiv:2204.03458} extended standard image diffusion architectures to the temporal domain, introducing conditional sampling for spatial and temporal video extension. Building upon this, MagicVideo (\texttt{arXiv:2211.11018}) optimized efficiency through a 3D U-Net design, a VideoVAE auto-encoder, and directed temporal attention modules. While MagicVideo offers a streamlined framework, it has been critiqued by \texttt{arXiv:2302.07685}, suggesting room for further refinement in temporal dependency modeling. Parallel to these structural developments, other models sought to address temporal length and continuity; Phenaki (\texttt{arXiv:2210.02399}) utilized a causal-attention-based video tokenizer and a bidirectional masked transformer to generate variable-length videos from open-domain textual descriptions. Collectively, these works represent a shift from the constrained, GAN-driven scene modeling of the mid-2010s to the highly flexible, text-guided latent diffusion frameworks that define current video generation capabilities.

\subsection{Advancements in Text-to-Video and Motion Control}

The evolution of text-to-video (T2V) generation has transitioned from leveraging pre-trained text-to-image (T2I) priors to developing specialized spatio-temporal architectures. Early methodologies like Make-A-Video bypass the need for paired text-video data by combining T2I models with spatiotemporally factorized pseudo-3D layers trained on unsupervised video [\texttt{arXiv:2209.14792}]. While this approach achieves high human preference (up to 84.38\% over VDM) and strong zero-shot performance on MSR-VTT (FID 13.17), it suffers from a fundamental limitation: it cannot learn text-to-motion associations for actions that are only inferable through video, such as "waving left-to-right" [\texttt{arXiv:2209.14792}]. In contrast, models like ModelScopeT2V utilize a larger 1.7 billion parameter architecture with dedicated spatio-temporal blocks to improve movement transitions [\texttt{arXiv:2308.06571}], though they have been critiqued by subsequent high-quality open models [\texttt{arXiv:2310.19512}].

A significant research branch focuses on efficient latent space manipulation to enable long or high-resolution synthesis. Latent Video Diffusion Models (LVDM) utilize a low-dimensional 3D latent space and hierarchical diffusion to generate videos exceeding 1000 frames [\texttt{arXiv:2211.13221}]. This concept of latent factorization is further refined by Projected Latent Video Diffusion Models (PVDM), which project videos into 2D-shaped latent vectors via an autoencoder to facilitate arbitrary length synthesis, achieving an FVD score of 639.7 on the UCF-101 benchmark [\texttt{arXiv:2302.07685}]. Similarly, Align Your Latents extends Latent Diffusion Models (LDMs) by introducing a temporal dimension to the latent space for high-resolution synthesis [\texttt{arXiv:2304.08818}]. Stable Video Diffusion advances this via a three-stage training strategy—comprising T2I pretraining, video pretraining, and high-quality finetuning—to scale these models to large datasets [\texttt{arXiv:2311.15127}].

Motion control and compositional synthesis represent the most recent frontier. While Tune-A-Video offers one-shot tuning for continuous motion using tailored spatio-temporal attention [\texttt{arXiv:2212.11565}], AnimateDiff provides a more flexible "plug-and-play" solution by learning transferable motion priors and utilizing MotionLoRA for lightweight fine-tuning [\texttt{arXiv:2307.04725}]. For more complex requirements, VideoComposer introduces a Spatio-Temporal Condition (STC) encoder, allowing for compositional synthesis using textual, spatial (sketch), and temporal (reference video) conditions [\texttt{arXiv:2306.02018}]. This level of control is complemented by the ControlNet architecture, which utilizes "zero convolutions" to add spatial conditioning—such as depth or human pose—to pretrained models [\texttt{arXiv:2302.05543}]. Finally, Text2Video-Zero offers a training-free alternative by reprogramming frame-level self-attention through cross-frame attention on the first frame, demonstrating that T2I models can function as zero-shot video generators [\texttt{arXiv:2303.13439}].

\subsection{Temporal Consistency and Video Editing Techniques}

Maintaining temporal consistency while enabling precise semantic manipulation represents a central challenge in video diffusion. Research in this domain diverges significantly based on whether the approach is training-free (zero-shot) or requires model finetuning.

Zero-shot methodologies focus on repurposing pre-trained image diffusion models for video tasks. Text2Video-Zero achieves this by enriching latent codes with motion dynamics and reprogramming frame-level self-attention through cross-frame attention on the first frame [\texttt{arXiv:2303.13439}]. While effective for generation, editing-specific zero-shot methods offer more granular control. Pix2Video utilizes a two-step process where edits are first applied to an anchor frame and then propagated via self-attention feature injection and latent code adjustment [\texttt{arXiv:2303.12688}]. In contrast, FateZero addresses the limitations of semantic leakage and structural preservation by fusing intermediate structural and motion attention maps and introducing spatial-temporal attention to the denoising UNet [\texttt{arXiv:2303.09535}]. Similarly, MotionFlow provides a training-free framework for motion transfer by utilizing cross-attention maps to manipulate spatial and temporal dynamics [\texttt{arXiv:2412.05275}]. For super-resolution tasks, Upscale-A-Video maintains temporal consistency through a combination of local temporal layers in the U-Net/VAE-Decoder and a global flow-guided recurrent latent propagation module [\texttt{arXiv:2312.06640}].

Conversely, training-based approaches seek to fundamentally restructure how models handle spatio-temporal dependencies. HiGen introduces a hierarchical decoupling strategy, separating spatial and temporal factors at both the structural (reasoning steps) and content (motion and appearance) levels to improve semantic accuracy and motion stability [\texttt{arXiv:2312.04483}]. VideoTetris advances compositional generation by employing spatio-temporal compositional diffusion and a reference frame attention mechanism [\texttt{arXiv:2406.04277}]. Addressing data quality issues, VideoCrafter2 proposes a training scheme that finetunes spatial modules with high-quality images to shift the distribution toward higher quality without inducing motion degradation [\texttt{arXiv:2401.09047}]. Finally, the Vectorized Timestep Approach (FVDM) redefines temporal modeling by using a vectorized timestep variable (VTV) to allow for independent frame noise schedules, enabling more fine-grained temporal dependencies [\texttt{arXiv:2410.03160}].

The literature reveals a clear tension between the flexibility of zero-shot methods and the stability of trained models. Zero-shot methods like FateZero [\texttt{arXiv:2303.09535}] and Text2Video-Zero [\texttt{arXiv:2303.13439}] offer immediate utility without retraining but may struggle with complex motion compared to hierarchical or decoupled architectures like HiGen [\texttt{arXiv:2312.04483}]. While training-free methods like MotionFlow [\texttt{arXiv:2412.05275}] provide efficient motion transfer, the structural improvements offered by specialized training schemes, such as the distribution shifting in VideoCrafter2 [\texttt{arXiv:2401.09047}] or the frame-aware scheduling in FVDM [\texttt{arXiv:2410.03160}], appear necessary to overcome the inherent limitations of standard image-to-video propagation.

\subsection{Efficient Architectures and Long-Term Generation}

The computational burden of video diffusion models stems from the multiplicative costs of spatial resolution, temporal length, and iterative denoising [\texttt{arXiv:2604.15911}]. Current research addresses these challenges through two primary trajectories: architectural innovations designed to replace quadratic-cost attention with more efficient mechanisms, and training-free strategies for extending generation duration.

Architectural advancements focus on optimizing temporal feature extraction. One approach leverages bidirectional State-Space Models (SSMs) to replace standard attention layers, enabling linear-time temporal extraction for long-term generation [\texttt{arXiv:2403.07711}]. Similarly, FrameDiT introduces "Matrix Attention," which processes entire frames as matrices using Frobenius inner products to preserve global spatio-temporal structures [\texttt{arXiv:2603.09721}]. While FrameDiT-H achieves significant efficiency and a 39\% FVD improvement over Latte on FaceForensics, its global-only variant (FrameDiT-G) struggles with temporal coherence, suggesting that matrix-level operations cannot yet fully replace local motion priors [\texttt{arXiv:2603.09721}]. Other methods seek parameter efficiency through motion-adaptive mechanisms; for instance, motion-adaptive temporal attention adjusts receptive fields—using local windows for high motion and global windows for low motion—achieving high efficiency with only 25.8M parameters [\texttt{arXiv:2603.17398}]. However, this lightweight approach suffers from lower generative quality (FID/FVD) and limited output scales compared to larger models [\texttt{arXiv:2603.17398}].

Efficiency is also pursued through latent decomposition and compression. The Content-Frame Motion-Latent Decomposition (CMD) model utilizes an autoencoder to separate video into a static content frame and a low-dimensional motion latent, achieving an FVD of 212.7 on WebVid-10M [\texttt{arXiv:2403.14148}]. CogVideoX further optimizes this via a 3D VAE for spatio-temporal compression and an expert transformer for text-video alignment [\texttt{arXiv:2408.06072}]. To optimize inference speed, VGDFR introduces a dynamic frame rate scheduler that adjusts latent space elements based on motion frequency, yielding up to a 3x speedup [\texttt{arXiv:2504.12259}].

A distinct category of research focuses on training-free methods to achieve long-duration or high-frame-rate video. GLC-Diffusion employs global-local collaborative denoising and noise reinitialization to maintain spatiotemporal consistency [\texttt{arXiv:2501.05484}], while VideoMerge extends generation by merging short clips from pretrained models [\texttt{arXiv:2503.09926}]. For high-frame-rate requirements, DiffuseSlide utilizes noise re-injection and sliding window latent denoising [\texttt{arXiv:2506.01454}]. Despite these advancements, a critical gap remains: current long-horizon generation relies heavily on "stability engineering" and "staged training" rather than unified algorithmic solutions or inherently stable objectives [\texttt{arXiv:2604.15911}].

\subsection{Current Frontiers: Survey and Advanced Representations}

The landscape of video diffusion models is characterized by a transition from broad architectural taxonomies toward highly specialized representations and modality-specific frameworks. Comprehensive surveys, such as \texttt{arXiv:2405.03150}, provide foundational overviews of architectural design and temporal dynamics modeling across diverse input modalities, including text, images, video, and audio [\texttt{arXiv:2405.03150}]. This is complemented by \texttt{arXiv:2504.16081}, which expands the scope to include optimization strategies, training engineering, and industry-specific solutions, with a heavy emphasis on computational efficiency and motion consistency [\texttt{arXiv:2504.16081}].

A significant shift in the literature is the treatment of image-to-video (I2V) generation as a distinct research domain rather than a mere sub-topic of text-to-video. \texttt{arXiv:2605.17248} establishes I2V as a standalone subject, introducing a unique two-level taxonomy and distilling four core technical design pillars: condition encoding, temporal modeling, noise prior design, and spatial-temporal upsampling [\texttt{arXiv:2605.17248}]. While this survey offers an exhaustive categorization of conditioning modalities—ranging from camera and motion to audio and multi-condition inputs—and provides a robust comparison of evaluation metrics, such as the VBench-I2V score, it possesses notable limitations [\texttt{arXiv:2605.17248}]. Specifically, the work lacks original empirical results and, while identifying a critical "reproducibility gap" caused by the rise of proprietary datasets, fails to propose a technical solution to this challenge [\texttt{arXiv:2605.17248}].

Moving beyond taxonomies, recent methodologies focus on enhancing the stability and awareness of internal representations. To address temporal coherence in text-to-video generation, RepVideo (\texttt{arXiv:2501.08994}) introduces a framework that accumulates features from neighboring layers to create enriched, stable semantic representations for attention mechanisms [\texttt{arXiv:2501.08994}]. In contrast, the Moaw framework (\texttt{arXiv:2601.12761}) pursues motion awareness through a different paradigm: shifting the diffusion modality from traditional image-to-video generation toward video-to-dense-tracking to enable zero-shot motion transfer [\texttt{arXiv:2601.12761}]. These methodological advancements are grounded in a deeper empirical understanding of latent structures; for instance, \texttt{arXiv:2502.07001} provides a systematic comparison between video and image diffusion architectures, specifically analyzing how temporal information influences latent representations for downstream visual understanding tasks [\texttt{arXiv:2502.07001}]. Collectively, these works demonstrate a field evolving from general generative modeling toward precise, controlled manipulation of motion and semantic stability.

\subsection{Tensions and Open Problems}

A primary tension in the literature concerns the architectural strategy for handling temporal dynamics: whether models should prioritize explicit temporal modeling or leverage high-fidelity spatial priors. Some researchers argue that video generation is best achieved by extending powerful image-based frameworks, as seen in the way \texttt{arXiv:2302.01329} utilizes the foundational capabilities of \texttt{arXiv:2006.11239}, \texttt{arXiv:2204.06125}, and \texttt{arXiv:2112.10752} to facilitate complex video editing. Conversely, others suggest that standard spatial models lack the necessary temporal coherence for long-form content. This is evidenced by the critique in \texttt{arXiv:2302.07685}, which argues that models like \texttt{arXiv:2209.14792} and \texttt{arXiv:2211.13221} struggle with structural consistency in projected latent spaces, suggesting that a more specialized approach to temporal dimensionality is required to prevent visual drift.

Another significant debate centers on the trade-off between training complexity and generative fidelity. One school of thought emphasizes the necessity of large-scale, text-video paired datasets to achieve high-quality results, a position implicitly critiqued by \texttt{arXiv:2209.14792}, which seeks to bypass such dependencies. However, more recent work suggests that simply increasing data or model scale is insufficient for nuanced control. For instance, \texttt{arXiv:2309.17444} argues that standard parameter-efficient methods like \texttt{arXiv:2106.09685} fail to provide the deep semantic grounding required for complex video synthesis, necessitating LLM-integrated architectures. This tension is further complicated by the pursuit of efficiency; while \texttt{arXiv:2312.04483} advocates for hierarchical spatio-temporal decoupling to improve upon the high-resolution capabilities of \texttt{arXiv:2307.01952}, others find that even these optimizations struggle to match the raw quality of massive-scale models like \texttt{arXiv:2310.19512}, which critiques the technical limitations found in the ModelScope framework of \texttt{arXiv:2308.06571}.

Finally, there is a growing disagreement regarding the metrics used to evaluate "success" in generative video. Traditional evaluation frameworks, often rooted in GAN-based stability or early video metrics like those in \texttt{arXiv:1706.08500} or \texttt{arXiv:1812.01717}, are increasingly viewed as inadequate for modern diffusion-based synthesis. The work in \texttt{arXiv:2311.17982} highlights this gap, arguing that existing benchmarks fail to capture the nuanced temporal and semantic qualities of modern video, thereby necessitating more comprehensive suites like VBench. This shift reflects a broader move away from simple distribution matching toward a more holistic assessment of motion realism, semantic alignment, and temporal consistency.

\subsection{Outlook}

The field is currently transitioning from "generation-at-all-costs" toward a paradigm of "controllable efficiency." As we move beyond the foundational breakthroughs of 2024, the frontier is defined by the integration of motion-aware architectures—such as the motion-adaptive attention mechanisms discussed in \texttt{arXiv:2603.17398}—and the development of training-free methods that allow for high-frame-rate synthesis without massive computational overhead (\texttt{arXiv:2506.01454}). The next generation of models will likely move away from fixed-rate latent representations toward dynamic, frame-rate-agnostic diffusion (\texttt{arXiv:2504.12259}), ultimately aiming to bridge the gap between static image aesthetics and the fluid, physics-compliant motion required for true cinematic synthesis.

\section*{References}
\begin{enumerate}\small
  \item UCF101: A Dataset of 101 Human Actions Classes From Videos in The Wild (2012). \texttt{arXiv:1212.0402}
  \item Auto-Encoding Variational Bayes (2013). \texttt{arXiv:1312.6114}
  \item Generative Adversarial Networks (2014). \texttt{arXiv:1406.2661}
  \item Deep Unsupervised Learning using Nonequilibrium Thermodynamics (2015). \texttt{arXiv:1503.03585}
  \item U-Net: Convolutional Networks for Biomedical Image Segmentation (2015). \texttt{arXiv:1505.04597}
  \item Generating Videos with Scene Dynamics (2016). \texttt{arXiv:1609.02612}
  \item Temporal Generative Adversarial Nets with Singular Value Clipping (2016). \texttt{arXiv:1611.06624}
  \item Attention Is All You Need (2017). \texttt{arXiv:1706.03762}
  \item GANs Trained by a Two Time-Scale Update Rule Converge to a Local Nash
  Equilibrium (2017). \texttt{arXiv:1706.08500}
  \item Decoupled Weight Decay Regularization (2017). \texttt{arXiv:1711.05101}
  \item Neural Discrete Representation Learning (2017). \texttt{arXiv:1711.00937}
  \item MoCoGAN: Decomposing Motion and Content for Video Generation (2017). \texttt{arXiv:1707.04993}
  \item Quo Vadis, Action Recognition? A New Model and the Kinetics Dataset (2017). \texttt{arXiv:1705.07750}
  \item The 2017 DAVIS Challenge on Video Object Segmentation (2017). \texttt{arXiv:1704.00675}
  \item A Style-Based Generator Architecture for Generative Adversarial Networks (2018). \texttt{arXiv:1812.04948}
  \item The Unreasonable Effectiveness of Deep Features as a Perceptual Metric (2018). \texttt{arXiv:1801.03924}
  \item Towards Accurate Generative Models of Video: A New Metric \& Challenges (2018). \texttt{arXiv:1812.01717}
  \item Video-to-Video Synthesis (2018). \texttt{arXiv:1808.06601}
  \item Train Sparsely, Generate Densely: Memory-efficient Unsupervised Training
  of High-resolution Temporal GAN (2018). \texttt{arXiv:1811.09245}
  \item Generative Modeling by Estimating Gradients of the Data Distribution (2019). \texttt{arXiv:1907.05600}
  \item Analyzing and Improving the Image Quality of StyleGAN (2019). \texttt{arXiv:1912.04958}
  \item Exploring the Limits of Transfer Learning with a Unified Text-to-Text
  Transformer (2019). \texttt{arXiv:1910.10683}
  \item Adversarial Video Generation on Complex Datasets (2019). \texttt{arXiv:1907.06571}
  \item Denoising Diffusion Probabilistic Models (2020). \texttt{arXiv:2006.11239}
  \item Denoising Diffusion Implicit Models (2020). \texttt{arXiv:2010.02502}
  \item Score-Based Generative Modeling through Stochastic Differential
  Equations (2020). \texttt{arXiv:2011.13456}
  \item An Image is Worth 16x16 Words: Transformers for Image Recognition at
  Scale (2020). \texttt{arXiv:2010.11929}
  \item Taming Transformers for High-Resolution Image Synthesis (2020). \texttt{arXiv:2012.09841}
  \item RAFT: Recurrent All-Pairs Field Transforms for Optical Flow (2020). \texttt{arXiv:2003.12039}
  \item First Order Motion Model for Image Animation (2020). \texttt{arXiv:2003.00196}
  \item Learning Transferable Visual Models From Natural Language Supervision (2021). \texttt{arXiv:2103.00020}
  \item High-Resolution Image Synthesis with Latent Diffusion Models (2021). \texttt{arXiv:2112.10752}
  \item Diffusion Models Beat GANs on Image Synthesis (2021). \texttt{arXiv:2105.05233}
  \item GLIDE: Towards Photorealistic Image Generation and Editing with
  Text-Guided Diffusion Models (2021). \texttt{arXiv:2112.10741}
  \item Zero-Shot Text-to-Image Generation (2021). \texttt{arXiv:2102.12092}
  \item Improved Denoising Diffusion Probabilistic Models (2021). \texttt{arXiv:2102.09672}
  \item Cascaded Diffusion Models for High Fidelity Image Generation (2021). \texttt{arXiv:2106.15282}
  \item VideoGPT: Video Generation using VQ-VAE and Transformers (2021). \texttt{arXiv:2104.10157}
  \item Frozen in Time: A Joint Video and Image Encoder for End-to-End Retrieval (2021). \texttt{arXiv:2104.00650}
  \item SDEdit: Guided Image Synthesis and Editing with Stochastic Differential
  Equations (2021). \texttt{arXiv:2108.01073}
  \item N\textbackslash{}"UWA: Visual Synthesis Pre-training for Neural visUal World creAtion (2021). \texttt{arXiv:2111.12417}
  \item GODIVA: Generating Open-DomaIn Videos from nAtural Descriptions (2021). \texttt{arXiv:2104.14806}
  \item A Good Image Generator Is What You Need for High-Resolution Video
  Synthesis (2021). \texttt{arXiv:2104.15069}
  \item LoRA: Low-Rank Adaptation of Large Language Models (2021). \texttt{arXiv:2106.09685}
  \item StyleGAN-V: A Continuous Video Generator with the Price, Image Quality
  and Perks of StyleGAN2 (2021). \texttt{arXiv:2112.14683}
  \item Make-A-Video: Text-to-Video Generation without Text-Video Data (2022). \texttt{arXiv:2209.14792}
  \item Hierarchical Text-Conditional Image Generation with CLIP Latents (2022). \texttt{arXiv:2204.06125}
  \item Photorealistic Text-to-Image Diffusion Models with Deep Language
  Understanding (2022). \texttt{arXiv:2205.11487}
  \item Video Diffusion Models (2022). \texttt{arXiv:2204.03458}
  \item Classifier-Free Diffusion Guidance (2022). \texttt{arXiv:2207.12598}
  \item Imagen Video: High Definition Video Generation with Diffusion Models (2022). \texttt{arXiv:2210.02303}
  \item CogVideo: Large-scale Pretraining for Text-to-Video Generation via
  Transformers (2022). \texttt{arXiv:2205.15868}
  \item Tune-A-Video: One-Shot Tuning of Image Diffusion Models for
  Text-to-Video Generation (2022). \texttt{arXiv:2212.11565}
  \item MagicVideo: Efficient Video Generation With Latent Diffusion Models (2022). \texttt{arXiv:2211.11018}
  \item DreamBooth: Fine Tuning Text-to-Image Diffusion Models for
  Subject-Driven Generation (2022). \texttt{arXiv:2208.12242}
  \item Prompt-to-Prompt Image Editing with Cross Attention Control (2022). \texttt{arXiv:2208.01626}
  \item Latent Video Diffusion Models for High-Fidelity Long Video Generation (2022). \texttt{arXiv:2211.13221}
  \item Scalable Diffusion Models with Transformers (2022). \texttt{arXiv:2212.09748}
  \item Long Video Generation with Time-Agnostic VQGAN and Time-Sensitive
  Transformer (2022). \texttt{arXiv:2204.03638}
  \item Generating Videos with Dynamics-aware Implicit Generative Adversarial
  Networks (2022). \texttt{arXiv:2202.10571}
  \item Progressive Distillation for Fast Sampling of Diffusion Models (2022). \texttt{arXiv:2202.00512}
  \item Flexible Diffusion Modeling of Long Videos (2022). \texttt{arXiv:2205.11495}
  \item An Image is Worth One Word: Personalizing Text-to-Image Generation using
  Textual Inversion (2022). \texttt{arXiv:2208.01618}
  \item Text2LIVE: Text-Driven Layered Image and Video Editing (2022). \texttt{arXiv:2204.02491}
  \item LAION-5B: An open large-scale dataset for training next generation
  image-text models (2022). \texttt{arXiv:2210.08402}
  \item Elucidating the Design Space of Diffusion-Based Generative Models (2022). \texttt{arXiv:2206.00364}
  \item Phenaki: Variable Length Video Generation From Open Domain Textual
  Description (2022). \texttt{arXiv:2210.02399}
  \item MCVD: Masked Conditional Video Diffusion for Prediction, Generation, and
  Interpolation (2022). \texttt{arXiv:2205.09853}
  \item eDiff-I: Text-to-Image Diffusion Models with an Ensemble of Expert
  Denoisers (2022). \texttt{arXiv:2211.01324}
  \item MoVideo: Motion-Aware Video Generation with Diffusion Models (2023). \texttt{arXiv:2311.11325}
  \item VideoGen: A Reference-Guided Latent Diffusion Approach for High
  Definition Text-to-Video Generation (2023). \texttt{arXiv:2309.00398}
  \item Stable Video Diffusion: Scaling Latent Video Diffusion Models to Large
  Datasets (2023). \texttt{arXiv:2311.15127}
  \item Hierarchical Spatio-temporal Decoupling for Text-to-Video Generation (2023). \texttt{arXiv:2312.04483}
  \item Video Probabilistic Diffusion Models in Projected Latent Space (2023). \texttt{arXiv:2302.07685}
  \item Align your Latents: High-Resolution Video Synthesis with Latent
  Diffusion Models (2023). \texttt{arXiv:2304.08818}
  \item LaMD: Latent Motion Diffusion for Image-Conditional Video Generation (2023). \texttt{arXiv:2304.11603}
  \item Upscale-A-Video: Temporal-Consistent Diffusion Model for Real-World
  Video Super-Resolution (2023). \texttt{arXiv:2312.06640}
  \item LLM-grounded Video Diffusion Models (2023). \texttt{arXiv:2309.17444}
  \item Adding Conditional Control to Text-to-Image Diffusion Models (2023). \texttt{arXiv:2302.05543}
  \item AnimateDiff: Animate Your Personalized Text-to-Image Diffusion Models
  without Specific Tuning (2023). \texttt{arXiv:2307.04725}
  \item Structure and Content-Guided Video Synthesis with Diffusion Models (2023). \texttt{arXiv:2302.03011}
  \item ModelScope Text-to-Video Technical Report (2023). \texttt{arXiv:2308.06571}
  \item VideoComposer: Compositional Video Synthesis with Motion Controllability (2023). \texttt{arXiv:2306.02018}
  \item LAVIE: High-Quality Video Generation with Cascaded Latent Diffusion
  Models (2023). \texttt{arXiv:2309.15103}
  \item Text2Video-Zero: Text-to-Image Diffusion Models are Zero-Shot Video
  Generators (2023). \texttt{arXiv:2303.13439}
  \item VBench: Comprehensive Benchmark Suite for Video Generative Models (2023). \texttt{arXiv:2311.17982}
  \item VideoCrafter1: Open Diffusion Models for High-Quality Video Generation (2023). \texttt{arXiv:2310.19512}
  \item T2I-Adapter: Learning Adapters to Dig out More Controllable Ability for
  Text-to-Image Diffusion Models (2023). \texttt{arXiv:2302.08453}
  \item VideoFusion: Decomposed Diffusion Models for High-Quality Video
  Generation (2023). \texttt{arXiv:2303.08320}
  \item Preserve Your Own Correlation: A Noise Prior for Video Diffusion Models (2023). \texttt{arXiv:2305.10474}
  \item Dreamix: Video Diffusion Models are General Video Editors (2023). \texttt{arXiv:2302.01329}
  \item SDXL: Improving Latent Diffusion Models for High-Resolution Image
  Synthesis (2023). \texttt{arXiv:2307.01952}
  \item FateZero: Fusing Attentions for Zero-shot Text-based Video Editing (2023). \texttt{arXiv:2303.09535}
  \item Pix2Video: Video Editing using Image Diffusion (2023). \texttt{arXiv:2303.12688}
  \item Video Diffusion Models: A Survey (2024). \texttt{arXiv:2405.03150}
  \item VideoTetris: Towards Compositional Text-to-Video Generation (2024). \texttt{arXiv:2406.04277}
  \item SSM Meets Video Diffusion Models: Efficient Long-Term Video Generation
  with Structured State Spaces (2024). \texttt{arXiv:2403.07711}
  \item Efficient Video Diffusion Models via Content-Frame Motion-Latent
  Decomposition (2024). \texttt{arXiv:2403.14148}
  \item VideoCrafter2: Overcoming Data Limitations for High-Quality Video
  Diffusion Models (2024). \texttt{arXiv:2401.09047}
  \item Lumiere: A Space-Time Diffusion Model for Video Generation (2024). \texttt{arXiv:2401.12945}
  \item Redefining Temporal Modeling in Video Diffusion: The Vectorized Timestep
  Approach (2024). \texttt{arXiv:2410.03160}
  \item MotionFlow: Attention-Driven Motion Transfer in Video Diffusion Models (2024). \texttt{arXiv:2412.05275}
  \item Grid Diffusion Models for Text-to-Video Generation (2024). \texttt{arXiv:2404.00234}
  \item Markov chains and mappings of distributions on compact spaces (2024). \texttt{arXiv:2404.01348}
  \item Markov chains and mappings of distributions on compact spaces II:
  Numerics and Conjectures (2024). \texttt{arXiv:2403.18153}
  \item CogVideoX: Text-to-Video Diffusion Models with An Expert Transformer (2024). \texttt{arXiv:2408.06072}
  \item Survey of Video Diffusion Models: Foundations, Implementations, and Applications (2025). \texttt{arXiv:2504.16081}
  \item Training-Free Motion-Guided Video Generation with Enhanced Temporal
  Consistency Using Motion Consistency Loss (2025). \texttt{arXiv:2501.07563}
  \item DiffuseSlide: Training-Free High Frame Rate Video Generation Diffusion (2025). \texttt{arXiv:2506.01454}
  \item From Image to Video: An Empirical Study of Diffusion Representations (2025). \texttt{arXiv:2502.07001}
  \item RepVideo: Rethinking Cross-Layer Representation for Video Generation (2025). \texttt{arXiv:2501.08994}
  \item VideoMerge: Towards Training-free Long Video Generation (2025). \texttt{arXiv:2503.09926}
  \item VGDFR: Diffusion-based Video Generation with Dynamic Latent Frame Rate (2025). \texttt{arXiv:2504.12259}
  \item Tuning-Free Long Video Generation via Global-Local Collaborative
  Diffusion (2025). \texttt{arXiv:2501.05484}
  \item Resource-Efficient Motion Control for Video Generation via Dynamic Mask
  Guidance (2025). \texttt{arXiv:2503.18386}
  \item Efficient Video Diffusion Models: Advancements and Challenges (2026). \texttt{arXiv:2604.15911}
  \item Image-to-Video Diffusion: From Foundations to Open Frontiers (2026). \texttt{arXiv:2605.17248}
  \item Motion-Adaptive Temporal Attention for Lightweight Video Generation with Stable Diffusion (2026). \texttt{arXiv:2603.17398}
  \item FrameDiT: Diffusion Transformer with Matrix Attention for Efficient Video Generation (2026). \texttt{arXiv:2603.09721}
  \item Moaw: Unleashing Motion Awareness for Video Diffusion Models (2026). \texttt{arXiv:2601.12761}
\end{enumerate}

\end{document}
